\subsection{Mechanism contrasts (RQ3)}

The study replays 8 matched-seed pairs per contrast. Each arm is a
separate history and fault episode, so the seed and schedule are matched while
the physical injections are independent. The counts below include every
planned pair; topology matches are reported explicitly. The term and role
snapshots are direct observations \cite{mongodb-replset-status}. The
interpretations are consistent with documented causal-session, read-concern,
and write-concern behavior
\cite{mongodb-causal-consistency,mongodb-read-concern,mongodb-write-concern};
they do not expose MongoDB's internal wait or replication state.

{\scriptsize
\setlength{\tabcolsep}{3pt}
\begin{longtable}{@{}p{1.7cm}p{2.5cm}p{7.8cm}p{2.4cm}@{}}
\caption{RQ3 matched-seed mechanism contrasts. Counts use all 8 pairs per contrast.}
\label{tab:rq3-mechanisms}\\
\toprule
Pair & Changed factor & Observed signatures & Topology / exemplar \\
\midrule
\endfirsthead
\caption[]{RQ3 matched-seed mechanism contrasts (continued).}\\
\toprule
Pair & Changed factor & Observed signatures & Topology / exemplar \\
\midrule
\endhead
\bottomrule
\endfoot
M1 (C5/C6) & causal: off/on & C5 stale v0/no bound: 8/8; C6 afterClusterTime present/matches W1 operationTime: 8/8/8/8; C6 reads issued: 8/8, successful: 0/8, not reached: 0/8; successful without a version: 0; nonstale concrete values: 0/0; stale among successful: 0/0. Starting primary and fault target matched: 8/8. & topology matched; selected \texttt{m1-r01} \\
M2 (C8/C5) & read concern: local/majority & C8 local v1: 8/8 (successful reads: 8/8); C5 majority v0: 8/8 (successful reads: 8/8); same route: 0/8; W2 acknowledgements C8/C5: 5/8/5/8; non-successes: 3/8/3/8; W2 on all final members C8/C5: 8/8/8/8; read version retained after convergence C8/C5: 0/8/8/8. Starting primary and fault target matched: 0/8. & topology differed; selected \texttt{m2-r01} \\
M3 (C3/C6) & write concern: w:1/majority & C3 w:1 first-write acknowledgements: 8/8; acknowledged W1 absent from every converged post-heal member: 8/8; C3 converged final states: 8/8; C6 majority acknowledgements: 0/8; C6 majority timeouts: 8/8; W1 later present on all converged members after NetworkTimeout: 8/8; C6 converged final states: 8/8. Starting primary and fault target matched: 4/8. & topology matched; selected \texttt{m3-r02} \\
\end{longtable}
}

The C3/C6 write comparison distinguishes acknowledgement from effect: a
\texttt{w:1} acknowledgement followed by W1 missing from every converged,
post-heal member is an observed rollback; a majority-write timeout remains
unresolved until the direct final-state observation. The C5/C6 causal comparison
records the command's causal time bound and client response separately, so a
timeout is not described as proof of a server-side wait. Read concern is
interpreted from the value and route actually recorded for each WFR read.

M2 classifications were C8: 5 VIOLATION/3 INDETERMINATE; C5: 5 PASS/3 INDETERMINATE. W2 appears on all final members in both arms, but C8's returned v1 is absent from the converged final versions in 8/8; C5's returned v0 remains in 8/8. Only 0/8 pairs matched both the initial primary and isolated member, so interpret this as a pattern under the recorded topologies rather than a fully topology-matched contrast.

\maybefigure[fig:rq3-causal-timeline]{submission/figures/rq3-causal-timeline.pdf}{Matched-seed C5/C6 RYW trace. The two arms use independent fault episodes; client command metadata, direct topology snapshots, and response outcomes are shown from the selected histories.}
